\section{Gauge / constraint dynamics, canonical GR, and selector-bearing quantum structure}
\label{app:gauge-qm}

\begin{fixedbox}
\textbf{Roadmap-only appendix.} This appendix uses \gaugeconstraintdoc{} as a theorem roadmap with split downstream payoff. The rigidity theorem proved there is not imported as a premise of Sections~1--9. Only its downstream classificatory payoff is used inside 42.
\end{fixedbox}

\subsection*{Scope and dependency declaration}
The source manuscript has two logically ordered layers: a theorem-bearing rigidity core for constrained Hamiltonian systems and a downstream classificatory payoff for canonical GR and selector-bearing quantum formalisms. In 42, only the second layer is imported. This appendix depends on the standing and role-discipline results of Sections~5--7 and on the no-backflow rule of Section~10.

\subsection*{Claim type}
This appendix is a theorem roadmap with two consequence-facing subparts. It does not inline the rigidity proof and does not promote the source manuscript into a live proof home for the foundational arc.

\subsection*{Canonical general relativity as a downstream bookkeeping projection}

\begin{proposition}[Canonical GR roadmap consequence]
Relative to the rigidity theorem of \gaugeconstraintdoc{}, canonical general relativity is to be read as a presentation-level realization of admissibility-derived constrained Hamiltonian structure rather than as the primitive source of that structure.
\end{proposition}

\begin{proof}
The source manuscript derives coisotropic standing, characteristic redescription, local constraint generators, first-class closure, and multiplier-type Hamiltonian structure from admissibility conditions and only afterward reads canonical GR through that structure. In that downstream reading, the ADM variables, constraint generators, and lapse/shift freedom become bookkeeping realizations of the admissible standing/redescription architecture rather than primitive explanatory inputs. The present appendix imports only that classificatory payoff.\cite{MaleyGaugeConstraintDynamics}
\end{proof}

\subsection*{Collapse, Everett weights, and selector-bearing quantum formalisms}

\begin{proposition}[Selector-bearing quantum consequence]
Within the same admissibility framework, collapse noise, Everettian branch weights, typicality measures, and related selector-bearing quantum formalisms do not acquire standing merely by being useful. Unless the selector structure is reified into physical state or proven quotient-invariantly forced, it remains meta-selection, added load, or bookkeeping.
\end{proposition}

\begin{proof}
This is the exact downstream diagnosis carried by the latter sections of \gaugeconstraintdoc{}. The source proves a collapse trilemma --- pure gauge, meta-selector, or reified noise --- and an Everett-weight trilemma --- pure bookkeeping, meta-selection, or reification --- and then strengthens the Everett point by exhibiting explicit non-invariance of branch decompositions under factorization and coarse-graining choices. In the 42 architecture, those results are not independent foundational theorems; they are the quantum-side payoff of treating selector-bearing structure as inadmissible on the original scope unless physically reified or quotient-forced.\cite{MaleyGaugeConstraintDynamics}
\end{proof}

\begin{reviewbox}
\textbf{Exact non-claims.} This appendix does not reproduce the rigidity derivation, does not claim to derive canonical GR from the capstone itself, and does not prove that all collapse or Everettian completions are impossible simpliciter. Its narrower claim is that, on the original scope, selector-bearing structures are bookkeeping or added load unless reified into state or forced as quotient invariants.\cite{MaleyGaugeConstraintDynamics}
\end{reviewbox}
